% !TEX program = xelatex
% Шаблон научной статьи НИТУ «МИСиС»
% Заполните блоки ниже и скомпилируйте: xelatex article.tex

\documentclass{misisarticle}

% --- Классификационные коды (оставьте пустым {}, если не требуются) ---
\grnti{47.01.07, 47.01.09}
\asjc{1707, 2207}
\oecd{1020, 1070}

% --- Название ---
\title{Название вашей научной статьи}

% --- Авторы: \misisauthor{ФИО}{аффилиация}{город, страна}{email} ---
\misisauthor{И.~О.~Фамилия}{кафедра инженерной кибернетики НИТУ «МИСиС»}{Москва, Россия}{student@edu.misis.ru}
\misisauthor{П.~П.~Научный}{кафедра инженерной кибернетики НИТУ «МИСиС»}{Москва, Россия}{supervisor@misis.ru}

% --- Аннотация и ключевые слова ---
\abstract{%
  Краткое изложение работы (5--10 предложений): постановка задачи, метод, основные результаты и выводы.
}

\keywords{%
  ключевое слово 1, ключевое слово 2, ключевое слово 3
}

\begin{document}

\maketitle

% ============================================================
\section{Введение}

Первый абзац раздела начинается с красной строки (отступ 0{,}51~см). Текст выравнивается по ширине, шрифт Times New Roman, 12~pt.

Второй абзац. Ссылки на литературу оформляйте так~\cite{example2024}. Для нескольких источников:~\cite{example2024, book2023}.

% ============================================================
\section{Основная часть}

\subsection{Подраздел}

Подзаголовки второго уровня набираются \textit{курсивом}.

Пример маркированного списка:
\begin{itemize}
  \item первый пункт;
  \item второй пункт;
  \item третий пункт.
\end{itemize}

Пример формулы:
\begin{equation}
  F_1 = 2 \cdot \frac{P \cdot R}{P + R},
  \label{eq:f1}
\end{equation}
где $P$~--- precision, $R$~--- recall.

% --- Рисунок ---
\begin{figure}[htbp]
  \centering
  % Замените заглушку на свой файл:
  % \includegraphics[width=0.75\linewidth]{figures/example}
  \fbox{\parbox[c][3cm][c]{0.75\linewidth}{\centering Placeholder для рисунка}}
  \caption{Пример иллюстрации. Замените заглушку на свой файл из папки figures.}
  \label{fig:example}
\end{figure}

% --- Таблица ---
\begin{table}[htbp]
  \caption{Пример таблицы с результатами}
  \label{tab:example}
  \centering
  \begin{tabular}{lcc}
    \toprule
    \textbf{Метод} & \textbf{Precision} & \textbf{Recall} \\
    \toprule
    Метод~A & 0.85 & 0.78 \\
    Метод~B & 0.91 & 0.83 \\
    \bottomrule
  \end{tabular}
\end{table}

% ============================================================
\section{Заключение}

Кратко сформулируйте основные выводы работы.

% ============================================================
% Список литературы — вариант 1: вручную
\begin{thebibliography}{9}

\bibitem{example2024}
Author A., Author B.
\newblock Title of the paper.
\newblock \textit{Journal Name}, vol.~1, no.~2, pp.~100--110, 2024.

\bibitem{book2023}
Author C.
\newblock \textit{Book Title}.
\newblock Publisher, City, 2023.

\end{thebibliography}

\end{document}
